
% Include necessary packages
\usepackage[nottoc,notbib]{tocbibind}  % biblatex adds the bibliography to the ToC itself

% References: biblatex keeps one reference segment per chapter, so each chapter
% can print its own reference list while the end of the book prints all of them.
% Numbering is global, so a source has the same number everywhere.
% backend=bibtex: biber's macOS binary fails on this machine; BibTeX works the same here.
\usepackage[backend=bibtex, style=numeric, sorting=none, refsegment=chapter, defernumbers=true]{biblatex}
\addbibresource{references.bib}
% Put \chapterreferences at the end of a chapter to list the sources cited in it.
% A chapter with no citations prints nothing.
\newcommand{\chapterreferences}{%
  \printbibliography[segment=\therefsegment, heading=subbibliography, title={References}]}
\usepackage{amsmath}
\usepackage{amsfonts}
\usepackage{amssymb}
\usepackage{graphicx}
\usepackage{float}
\usepackage{amsthm}
\usepackage{hyperref}
\usepackage{cleveref}
\usepackage{caption}
\usepackage{subcaption}
\usepackage{enumitem}
\usepackage{appendix}
\usepackage{titlesec}
\usepackage{setspace}
\usepackage{geometry}
\usepackage{pdflscape}
\usepackage{pdfpages}
\usepackage{wrapfig}
\usepackage{fancyhdr}
\usepackage{etoolbox}
\usepackage{mathrsfs}
\usepackage{tikz}
\usepackage{tikz-cd}
\usepackage{pgfplots}
\pgfplotsset{compat=1.18}
\usepackage{pgfplotstable}
\usepackage{booktabs}
\usepackage{array}
\usepackage{multirow}
\usepackage{longtable}
\usepackage{listings}
\usepackage{color}
\usepackage{colortbl}
\usepackage{braket}
\usepackage{quantikz}
\usepackage{circuitikz}
\usepackage{algorithm}
\usepackage{algpseudocode}
\usepackage{tcolorbox}

% Define colors
\definecolor{codegreen}{rgb}{0,0.6,0}
\definecolor{codegray}{rgb}{0.5,0.5,0.5}
\definecolor{codepurple}{rgb}{0.58,0,0.82}
\definecolor{backcolour}{rgb}{0.95,0.95,0.92}
\definecolor{darkgreen}{rgb}{0.0, 0.2, 0.13}
\definecolor{darkred}{rgb}{0.55, 0.0, 0.0}
\definecolor{darkblue}{rgb}{0.0, 0.0, 0.55}
\definecolor{darkpurple}{rgb}{0.5, 0.0, 0.5}
\definecolor{darkcyan}{rgb}{0.0, 0.55, 0.55}
\definecolor{darkgray}{rgb}{0.66, 0.66, 0.66}
\definecolor{lightgray}{rgb}{0.95, 0.95, 0.95}

% Define styles and environments
\lstdefinestyle{mystyle}{
    backgroundcolor=\color{backcolour},   
    commentstyle=\color{codegreen},
    keywordstyle=\color{magenta},
    numberstyle=\tiny\color{codegray},
    stringstyle=\color{codepurple},
    basicstyle=\footnotesize,
    breakatwhitespace=false,         
    breaklines=true,                 
    captionpos=b,                    
    keepspaces=true,                 
    numbers=left,                    
    numbersep=5pt,                  
    showspaces=false,                
    showstringspaces=false,
    showtabs=false,                  
    tabsize=2
}
\lstset{style=mystyle}

\newtcolorbox{importantnote}{
    colback=lightgray,
    colframe=darkgray,
    fonttitle=\bfseries,
    title=Important Note
}

\newtheorem{problem}{Problem}[section]
\newtheorem{theorem}{Theorem}[section]
\newtheorem{lemma}[theorem]{Lemma}
\newtheorem{corollary}[theorem]{Corollary}
\newtheorem{proposition}[theorem]{Proposition}
\newtheorem{axiom}{Axiom}[section]
\theoremstyle{definition}
\newtheorem{definition}{Definition}[section]
\theoremstyle{definition}
\newtheorem{example}{Example}[section]
\theoremstyle{remark}
\newtheorem*{remark}{Remark}

% Define \abstractname and \acknowledgementsname
\newcommand{\abstractname}{Abstract}
\newcommand{\acknowledgementsname}{Acknowledgements}

\newtcolorbox{note}{
    colback=lightgray,
    colframe=darkgray,
    fonttitle=\bfseries,
    title=Note
}
\newenvironment{abstract}{%
\clearpage
\null\vfill
\begin{center}%
    \bfseries \abstractname
\end{center}}%
{\vfill\null}

\newenvironment{Acknowledgements}{%
\clearpage
\null\vfill
\begin{center}%
    \bfseries \acknowledgementsname
\end{center}}%
{\vfill\null}

